\reportelement{Введение}
Цель работы --- подключаем к приложению PostgreSQL, учимся хранить данные вне контейнера и разбираемся с сетями Docker.

Работа выполнена в Ubuntu WSL с использованием Docker Desktop. Продолжен проект второй лабораторной работы: приложение FastAPI расположено в каталоге \textnormal{solutions/assist-py/}, изменения сохранены в ветке \textnormal{lab-3} личного форка.

Последние две цифры номера ИСУ --- 28. Порт PostgreSQL на хосте рассчитан как $5430+28=5458$, порт приложения --- как $8000+28=8028$. Персональные параметры приведены в таблице~\ref{tab:params}.

\begin{table}[H]
\caption{Персональные параметры}\label{tab:params}
\centering
\renewcommand{\arraystretch}{1.12}
\begin{tabularx}{\linewidth}{|>{\raggedright\arraybackslash}X|>{\raggedright\arraybackslash}p{.37\linewidth}|}
\hline
Параметр & Значение \\\hline
Номер ИСУ & 467228 \\\hline
Порт PostgreSQL на хосте & 5458 \\\hline
Порт приложения на хосте & 8028 \\\hline
Персональный префикс & rekun \\\hline
Учебный секретный ключ & secret28rekun \\\hline
Имя базы данных & rekun28 \\\hline
\end{tabularx}
\end{table}

Для базы данных используется контейнер \textnormal{rekun-postgres}, для приложения --- \textnormal{rekun-assistant-lab3}. Созданы именованный том \textnormal{rekun-pgdata} и пользовательская сеть \textnormal{rekun-backend}.

Репозиторий и версия приложения для третьей работы: \url{https://github.com/grettl/course-containerization-orchestration-application/tree/lab-3}.

\chapter{ЗАПУСК POSTGRESQL И ПРИЛОЖЕНИЯ}
\section{Создание тома и запуск базы данных}
\screenshotblock{1}{Создание и проверка тома}{Для хранения файлов PostgreSQL создан именованный том \textnormal{rekun-pgdata}. Его наличие подтверждено командой \textnormal{docker volume ls} (рисунок~\ref{fig:1}). Том существует отдельно от файловой системы контейнера.}{Команда docker volume create создаёт том, а docker volume ls показывает список томов. Docker управляет местом их хранения. При удалении контейнера именованный том остаётся, если его не удалить отдельно.}

\screenshotblock{2}{Контейнер PostgreSQL}{Контейнер \textnormal{rekun-postgres} запущен на образе \textnormal{postgres:16}. Том подключён к каталогу \textnormal{/var/lib/postgresql/data}. Порт хоста \textnormal{5458} сопоставлен с внутренним портом базы \textnormal{5432} (рисунок~\ref{fig:2}).}{В записи 5458:5432 слева указан порт компьютера, справа --- порт PostgreSQL внутри контейнера. Подключение тома к каталогу данных позволяет хранить файлы базы отдельно от контейнера. Статус Up означает, что контейнер запущен.}

\screenshotblock{3}{Готовность PostgreSQL к подключениям}{Готовность базы проверена утилитой \textnormal{pg\_isready}. Ответ \textnormal{accepting connections} подтверждает, что сервер принимает подключения (рисунок~\ref{fig:3}).}{Команда docker exec выполняет проверку в работающем контейнере. Утилита pg\_isready проверяет, принимает ли сервер подключения; параметры -U и -d задают пользователя и базу. Это более конкретная проверка готовности PostgreSQL, чем наличие статуса Up.}

\Needspace{.55\textheight}
\section{Подключение FastAPI к PostgreSQL}
В приложение добавлены зависимость \textnormal{psycopg2-binary} и модуль \textnormal{app/database.py}. При запуске создаётся таблица \textnormal{notes}, если её ещё нет. Обработчик \textnormal{POST /notes} записывает заметку в базу, а \textnormal{GET /notes} возвращает сохранённые записи. Параметры соединения передаются через переменные окружения.

\screenshotblock{4}{Проверка приложения и сохранения заметки}{После сборки образа \textnormal{rekun-assistant:lab3} приложение запущено с публикацией порта \textnormal{8028:5458}. Выполнены четыре запроса: к главной странице, к маршруту \textnormal{/health}, на создание заметки и на получение списка. Новая заметка «ЛР3: данные в БД» получила идентификатор 3. В списке отображаются все три сохранённые записи (рисунок~\ref{fig:4}).}{curl отправляет HTTP-запросы приложению. GET получает данные, POST создаёт запись; тело POST передаётся в формате JSON. Ответ healthy показывает работу маршрута проверки приложения, а успешные запись и чтение заметки подтверждают взаимодействие с базой.}

\chapter{ПРОВЕРКА СОХРАНЕНИЯ ДАННЫХ}
\screenshotblock{5}{Удаление контейнера приложения}{Для проверки сохранения данных контейнер \textnormal{rekun-assistant-lab3} остановлен и удалён. На рисунке~\ref{fig:5} показаны команды удаления при первоначальной проверке. Контейнер PostgreSQL и именованный том продолжили хранить данные.}{docker stop останавливает процесс контейнера, docker rm удаляет сам контейнер. Образ приложения остаётся и используется для нового запуска. Здесь удаляется только FastAPI: PostgreSQL продолжает работать, поэтому сохранённые в нём записи не теряются.}

\screenshotblock{6}{Заметка после пересоздания приложения}{После добавления трёх заметок контейнер приложения повторно пересоздан. Он подключён к сети \textnormal{rekun-backend} с прежними параметрами доступа к PostgreSQL. Запрос \textnormal{GET /notes} вернул все три записи. На рисунке~\ref{fig:6} видны дата и время проверки, логин, имя компьютера и список заметок.}{Персистентность означает сохранение данных независимо от времени жизни приложения. Если бы заметки оставались только в списке Python, новый процесс начал бы с пустого списка. Получение трёх прежних записей показывает, что приложение читает их из PostgreSQL.}

\chapter{ПОДКЛЮЧЕНИЕ КОНФИГУРАЦИИ}
В каталоге проекта расположен файл \textnormal{config/app.conf}. Он подключён через bind mount к пути \textnormal{/app/config/app.conf} в проверочном контейнере \textnormal{rekun-config-check}. Использован режим \textnormal{readonly}, который ограничивает доступ контейнера к файлу чтением.

\screenshotblock{7}{Конфигурация внутри контейнера}{Файл приведён к образцу задания: задано ограничение количества заметок, включён параметр погодного API, указаны \textnormal{maintenance\_mode=false} и владелец \textnormal{owner=rekun}. Содержимое прочитано командой \textnormal{docker exec} внутри контейнера (рисунок~\ref{fig:7}).}{При bind mount конкретный файл компьютера становится доступен по заданному пути контейнера. В отличие от именованного тома, здесь явно выбирается исходный путь на хосте. Команда cat выводит содержимое файла; owner указывает владельца учебной конфигурации.}

\screenshotblock{8}{Изменение конфигурации без пересборки}{На хосте значение \textnormal{maintenance\_mode} изменено с \textnormal{false} на \textnormal{true}. После перезапуска контейнера повторное чтение файла внутри него показало новое значение (рисунок~\ref{fig:8}). Docker-образ при этом не пересобирался.}{readonly запрещает изменение подключённого файла из контейнера, но позволяет менять его на хосте. docker restart перезапускает существующий контейнер. Файл подключён снаружи, поэтому для изменения его содержимого новая сборка образа не требуется.}

В этой реализации файл используется для проверки монтирования; настройки соединения FastAPI с базой передаются через переменные окружения.

\chapter{ПОЛЬЗОВАТЕЛЬСКАЯ СЕТЬ И ИЗОЛЯЦИЯ}
\section{Создание сети и подключение контейнеров}
\screenshotblock{9}{Создание сети rekun-backend}{Создана пользовательская сеть \textnormal{rekun-backend} с драйвером \textnormal{bridge}. Результат создания и наличие сети в списке показаны на рисунке~\ref{fig:9}.}{Пользовательская bridge-сеть обеспечивает взаимодействие подключённых контейнеров на одном Docker-хосте. Её имя задаётся при запуске параметром --network. Внутри такой сети Docker позволяет обращаться к другим контейнерам по их именам.}

\screenshotblock{10}{Контейнеры в пользовательской сети}{К сети подключены контейнеры базы данных и приложения. Команда \textnormal{docker network inspect} выводит имена \textnormal{rekun-postgres} и \textnormal{rekun-assistant-lab3} (рисунок~\ref{fig:10}).}{docker network inspect показывает подробные сведения о сети. В данном выводе выбраны только имена подключённых контейнеров. Наличие обоих имён подтверждает их участие в одной сети; реальный обмен данными проверяется отдельно запросом к приложению.}

\screenshotblock{11}{Доступ к заметке после настройки сети}{Приложение пересоздано с параметрами \textnormal{DB\_HOST=rekun-postgres} и \textnormal{DB\_PORT=5432}. Подключение выполняется по имени контейнера через внутреннюю сеть Docker. При первоначальной проверке запрос к API вернул созданную ранее заметку (рисунок~\ref{fig:11}).}{Внутри контейнера localhost обозначает его самого. Поэтому адрес базы задан именем rekun-postgres. Для связи между контейнерами используется внутренний порт PostgreSQL 5432; опубликованный порт хоста 5458 предназначен для подключения через компьютер.}

\Needspace{.45\textheight}
\section{Проверка связи между контейнерами}
\screenshotblock{12}{Успешный ping внутри сети}{Из временного контейнера в сети \textnormal{rekun-backend} выполнен \textnormal{ping rekun-postgres}. Имя разрешилось в адрес \textnormal{172.18.0.2}; получены ответы на оба пакета, потери составили 0\% (рисунок~\ref{fig:12}).}{ping проверяет доступность узла пакетами ICMP. Параметр -c 2 ограничивает проверку двумя пакетами. Вывод показывает разрешение имени в IP-адрес и получение двух ответов: 0 процентов потерь означает успешную проверку.}

\screenshotblock{13}{Неудачный ping из сети bridge}{Из временного контейнера в стандартной сети \textnormal{bridge} проверен тот же IP-адрес базы. Ответы не получены: потери составили 100\% (рисунок~\ref{fig:13}). Для сопоставления использован фактический адрес PostgreSQL, полученный через \textnormal{docker inspect}.}{Вторая проверка выполняется из другой Docker-сети по фактическому IP базы. Потери 100 процентов означают, что в этом опыте ICMP-ответы не пришли. Этот результат характеризует проверенный маршрут и протокол; он не проверяет доступ через опубликованный порт хоста.}

В выполненном опыте ICMP-связь с контейнером базы доступна внутри пользовательской сети и отсутствует из стандартной сети \textnormal{bridge}.

\chapter{ИТОГОВАЯ ПРОВЕРКА}
\screenshotblock{14}{Контейнеры, том и сеть Docker}{Выполнена финальная проверка с указанием даты, времени, логина и имени компьютера. Показаны запущенные контейнеры \textnormal{rekun-postgres}, \textnormal{rekun-assistant-lab3} и \textnormal{rekun-config-check}, именованный том \textnormal{rekun-pgdata} и пользовательская сеть \textnormal{rekun-backend} (рисунок~\ref{fig:14}).}{docker ps показывает работающие контейнеры, docker volume ls --- тома, docker network ls --- сети. Фильтр name=rekun- оставляет объекты с персональным префиксом. Дата, логин и хост позволяют связать итоговую проверку с окружением студента.}

В списке контейнеров видны публикации портов: \textnormal{8028:5458} для FastAPI и \textnormal{5458:5432} для PostgreSQL. Наличие тома и сети подтверждено отдельными командами Docker.

В репозитории сохранены исходный код, файл конфигурации и инструкция по воспроизведению работы. В README описаны создание тома и сети, запуск PostgreSQL и FastAPI, проверка заметок, bind mount и сетевого взаимодействия.

\reportelement{Заключение}
К приложению FastAPI подключена база PostgreSQL, заметки перенесены в таблицу базы данных. Проверено сохранение заметки после удаления и пересоздания контейнера приложения.

Для файлов PostgreSQL создан именованный том. Конфигурационный файл подключён через bind mount в режиме чтения. Проверено обновление файла внутри контейнера после изменения на хосте без пересборки образа. Контейнеры приложения и базы объединены в пользовательскую сеть; проверены связь внутри неё и отсутствие ICMP-ответов из стандартной сети \textnormal{bridge}.
